\section{The first unavailable split and risk on future cases}
\label{app:first-split}
Fix the forest and a retained support $S$. For one labeled row $(x,y)$,
write
\[
  H_s(x,y;S)=\frac{p}{s}\,\mathbb E_T
  \bigl[\ell(f_{T\cap S}(x),y)-\ell(f_T(x),y)\bigr],
  \qquad |T|=s ,
\]
where $T$ is conditionally uniform at size $s$ and independent of the row.
The forest averages its tree outputs with nonnegative weights summing to
one. For Brier loss set $L=2$. For regression assume that predictions and
responses lie in an interval of width $R$ and set $L=R^2$.

Follow each tree's full-input path on $x$ to its first split outside $S$.
Let $\pi(S,x)$ be the total weight of the trees with such a split, and let
$j$ be its feature in one of these trees. If $j\notin T$, the two traversals
stop together at that split or at an earlier one. Thus the tree can change
only if $j\in T$, an event of probability $s/p$. Trees without such a split
follow identical paths. If $\theta$ is the total weight of trees whose
outputs change, then $\mathbb E_T\theta\le(s/p)\pi(S,x)$.

Two regression forest outputs differ by at most $R\theta$, and their
squared errors differ by at most $2R^2\theta$. Two class probability vectors
$u,v$ differ by at most $2\theta$ in $\ell_1$, since the unchanged trees
contribute equally to both. For a one-hot $y$, the absolute Brier difference
$|\langle u-v,u+v-2y\rangle|$ is therefore at most $4\theta$. In both cases
the individual loss also lies in $[0,L]$, giving
\[
  |H_s(x,y;S)|\le L\min\{p/s,\,2\pi(S,x)\}.
\]
This controls the factor $p/s$ even at size one. For a nonuniform mask,
the path argument uses the actual conditional probability of acquiring
its first excluded feature. A uniform $s/p$ bound then needs additional
assumptions. Averaging probabilities or bounded scalar tree outputs is
essential to the constants above. A sigmoid applied to a sum of boosted
margins would require bounds on the individual margin changes and the link.

For a statistical consequence, condition on the forest and $K$ candidate
supports chosen independently of $n$ independent and identically distributed
acceptance rows. Assume future rows follow the same distribution. Put
$X_s=H_s-\epsilon\ell(f_\varnothing(x),y)$ and
$B_s(S)=L\min\{p/s,\,2\sup_x\pi(S,x)\}$. By the one-sided Hoeffding
inequality~\citep{hoeffding1963probability} and a union bound, with probability
at least $1-\delta$ every support and every positive size satisfy
\[
  \mathbb E X_s\ \le\ \overline X_s+
      \bigl(2B_s(S)+\epsilon L\bigr)
      \sqrt{\frac{\log(Kp/\delta)}{2n}} .
\]
Imposing nonpositivity of these upper bounds would control the population
quota for any support selected from the family. The experiments instead
impose the empirical quota. For $p=16$, $K=34$, $n=300$, $\delta=.05$,
$\epsilon=.02$, Brier loss and $\sup_x\pi=1$, the added margins are
$1.00$ at $s=1$ and $.503$ at $s=p$. A balanced binary null prediction has
Brier loss $.5$, giving an allowed excess of only $.01$ per available
fraction. Chronological changes in the row distribution also fall outside
the sampling assumptions.

\paragraph{Limit under unused inputs.}\label{app:padding}
For fixed $S$, write $G(t)=\sum_{U\not\subseteq S} w_U t^{|U|-1}$.
When each input is independently available with probability $t>0$,
$G(t)$ is the expected excess loss divided by $t$. Conditioning on the
available count shows that its Bernstein coefficients at degree $p-1$
are exactly $V_1^{(p)},\ldots,V_p^{(p)}$: the coefficient of vote $s$
in the mixture is $\binom{p-1}{s-1} t^{s-1}(1-t)^{p-s}$. Adding one
unused input therefore performs ordinary degree elevation. For $p>1$,
\[
0\ \le\ \max_s V_s^{(p)}-\max_{0\le t\le1}G(t)\
\le\ \sum_{U\not\subseteq S}|w_U|
          \min\!\left\{1,\frac{\binom{|U|-1}{2}}{p-1}\right\}.
\]
The lower bound follows from the mixture. For the upper bound, put
$t=(s-1)/(p-1)$ and $k=|U|-1$. The binomial ratio in $V_s$ is the
probability that $k$ draws without replacement from $p-1$ objects all
land among $s-1$ marked objects. With replacement the probability is
$t^k$. Couple the draws until the first repeated object; its probability
is at most $\binom{k}{2}/(p-1)$ by the union bound. Summing the resulting
errors proves the inequality. With the forest fixed, padding can only
decrease the maximum, and it converges to the worst independent-availability
rate. Fixed-size and independent-availability guarantees can therefore
differ even for exactly the same model and rows.
